\section{\texorpdfstring{Translating the definition into a Lean \texttt{structure}}{Translating the definition into a Lean structure}}\label{sec:07-groups:02-translating:1}

The definition names three pieces of data and three families of axioms. Lean has no primitive for ``data plus proof obligations'' other than \texttt{structure} itself, so the translation proceeds in two steps: record the data first, with nothing checked, then add one field per axiom.

\textbf{Step 1. Just the data.}

\begin{lstlisting}
structure GroupData (G : Type) where
  op : G (*$\rightarrow$*) G (*$\rightarrow$*) G
  id : G
  inv : G (*$\rightarrow$*) G
\end{lstlisting}

A \texttt{GroupData G} is an operation, an identity element, and an inverse function, with no axiom checked. It can still be nonsense: \texttt{op} could ignore both arguments.

\begin{mathreading}

\texttt{GroupData G} is the \emph{signature} of a group on the carrier \(G\), the raw data \((\,\cdot : G\times G \to G,\ e \in G,\
(-)^{-1} : G \to G\,)\) with no laws imposed, the product \[
\mathrm{GroupData}(G) = (G^{G\times G}) \times G \times (G^{G}).
\] It is a \emph{magma} signature, not yet a group; the axioms below cut out the genuine group structures as a subset.

\end{mathreading}

\textbf{Step 2. Add the axioms as extra fields.} A proof can be a field of a structure exactly as data can. One field per axiom, each of type ``a proposition that must hold for every element'':

\begin{lstlisting}
structure Group (G : Type) where
  op : G (*$\rightarrow$*) G (*$\rightarrow$*) G
  id : G
  inv : G (*$\rightarrow$*) G
  assoc : (*$\forall$*) a b c : G, op (op a b) c = op a (op b c)
  id_left : (*$\forall$*) a : G, op id a = a
  id_right : (*$\forall$*) a : G, op a id = a
  inv_left : (*$\forall$*) a : G, op (inv a) a = id
  inv_right : (*$\forall$*) a : G, op a (inv a) = id
\end{lstlisting}

\texttt{assoc} is a \emph{proof obligation}: whoever constructs a \texttt{Group G} must supply a term proving it for every \texttt{a b c}. The remaining four fields are the identity and inverse axioms, split left/right because commutativity has not been assumed. Collapsing them into one field each would silently assume it.

\begin{mathreading}

\texttt{Group G} is the type of \emph{group structures on the fixed carrier \(G\)}, a dependent tuple \[
\Big(\,\cdot,\ e,\ (-)^{-1},\ \underbrace{\alpha}_{\text{assoc}},\
\underbrace{\lambda_\ell,\lambda_r}_{\text{identity}},\
\underbrace{\iota_\ell,\iota_r}_{\text{inverse}}\,\Big),
\] where the last five components are proofs of \[
\forall a,b,c,\ (a\cdot b)\cdot c = a\cdot(b\cdot c);\quad
\forall a,\ e\cdot a = a;\ \ldots;\quad \forall a,\ a\cdot a^{-1}=e.
\] \texttt{Group G} is exactly the subset of \texttt{GroupData G} cut out by these five propositions, the \(\Sigma\)-type \(\sum_{d : \mathrm{GroupData}(G)}
\mathrm{Axioms}(d)\), the same shape as Chapter 4's ``witness together with a proof about it,'' here with a whole bundle of data as the witness. This matches ``a set with operations \emph{such that} the axioms hold''; the ``such that'' is a genuine subobject of the space of raw data.

\end{mathreading}

\begin{quote}
Read more. The actual \texttt{Group} in Mathlib (\texttt{Mathlib.Algebra.Group.Defs}) is a \texttt{class}, not the plain \texttt{structure} used in this book, inheriting from a chain of smaller classes (\texttt{Mul}, \texttt{One}, \texttt{Inv}, \texttt{Monoid}, \ldots) instead of listing all axioms in one place. See \hyperref[sec:14-next-steps:02-moving-to-mathlib:1]{Chapter 14} for the bridge between the two styles, and \hyperref[sec:06-rigor-check:01-structure-vs-class:1]{Chapter 6, Section 1} for why this book delays that mechanism.
\end{quote}

\begin{progcorner}

An ordinary Python class checks none of this:

\end{progcorner}

\begin{lstlisting}[language=Python,style=python]
class Group:
    def __init__(self, op, id, inv):
        self.op = op
        self.id = id
        self.inv = inv
\end{lstlisting}

\texttt{Group(op=lambda a, b: a -\allowbreak{} b, id=0, inv=lambda a: a)} type-checks and runs, and is not a group, since subtraction is not associative. The bug surfaces later, silently, wherever a theorem assuming associativity is applied to this \texttt{op} without re-checking the assumption. \texttt{GroupData} above has the identical defect. Values are bundled but nothing is checked, which is exactly why \texttt{Group} adds the five axiom fields: a term of type \texttt{Group G} cannot exist unless proofs of \texttt{assoc}, \texttt{id\_\allowbreak{}left}, \texttt{id\_\allowbreak{}right}, \texttt{inv\_\allowbreak{}left}, and \texttt{inv\_\allowbreak{}right} were actually supplied, so every theorem in the rest of the book taking a \texttt{Group G} gets those five facts for free, checked once, rather than trusting every caller to have built the underlying data correctly.
